\section*{Bab \ref{ch:b2:biomolecules} --- Biomolekul: Asam Amino, Gula, Lipid, dan Nukleotida}

\begin{solution}{exo:b2:biomolecules:1}
\begin{itemize}
  \item Di antara pH 3 dan 9, \omterm{def:b2:biomolecules:amino-acid}{zwitterion} \ce{H3N+-CH(CH3)-COO-}.
  \item Pada pH 1, kation \ce{H3N+-CH(CH3)-COOH}.
  \item Pada pH 12, anion \ce{H2N-CH(CH3)-COO-}.
\end{itemize}
\end{solution}

\begin{solution}{exo:b2:biomolecules:2}
Glisina: $(2.35 + 9.78)/2 \approx 6.1$. Alanina: $(2.35 + 9.87)/2 \approx 6.1$.
% ledger: pka:glycine1, pka:glycine2, pka:alanine1, pka:alanine2
\end{solution}

\begin{solution}{exo:b2:biomolecules:3}
\ce{H2N-CH(CH3)-CO-NH-CH2-COOH}: alanina menyediakan ujung-N, glisina ujung-C.
Gly--Ala, \ce{H2N-CH2-CO-NH-CH(CH3)-COOH}, adalah isomer konstitusi: gugus
amino bebasnya berada pada glisina.
\end{solution}

\begin{solution}{exo:b2:biomolecules:4}
\ce{NH2} di kiri: L. Prioritas \ce{NH2} > \ce{COOH} > \ce{CH2OH} > \ce{H}: L-serina adalah
(\textit{S}).
\end{solution}

\begin{solution}{exo:b2:biomolecules:5}
Gugus-gugus: amonium ujung-N (sekitar 9), \omterm{def:b2:biomolecules:amino-acid}{rantai samping} lisina (10.7),
\omterm{def:b2:biomolecules:amino-acid}{rantai samping} aspartat (3.9), gugus karboksilat ujung-C (sekitar 2). pH 1:
$+1 +1 + 0 + 0 = +2$. pH 7: $+1 +1 -1 -1 = 0$. pH 12:
$0 + 0 - 1 - 1 = -2$ (\omterm{def:b2:biomolecules:amino-acid}{rantai samping} lisina, 1.3 satuan di bawah pH 12,
sebagian besar netral).
\end{solution}

\begin{solution}{exo:b2:biomolecules:6}
Tionil klorida dan pemanasan juga akan menyerang gugus-gugus lain dan,
terutama, \omterm{def:b2:acyl-substitution:derivative}{asil klorida} sebuah asam amino yang sangat reaktif mudah
kehilangan konfigurasinya (hidrogen $\alpha$-nya menjadi asam): \omterm{def:b2:biomolecules:peptide}{peptida} itu
akan sebagian terasemisasi. \omterm{def:b2:biomolecules:peptide}{Pereaksi kopling} mengaktifkan asam pada suhu
kamar, dalam kondisi lunak.
\end{solution}

\begin{solution}{exo:b2:biomolecules:7}
O5 terikat pada C2: cincin C2, C3, C4, C5, dan O, lima atom (sebuah
furanosa). Dalam \omterm{def:b2:biomolecules:anomer}{proyeksi Haworth}, O cincin di belakang, C2 di kanan dengan
\ce{CH2OH} (C1) di bawah dan \ce{OH} di atas ($\beta$); \ce{OH} C3 di atas
(di kiri dalam proyeksi Fischer), \ce{OH} C4 di bawah, C5 membawa \ce{CH2OH}
(C6) di atas.
\end{solution}

\begin{solution}{exo:b2:biomolecules:8}
$x_\alpha = (80.2 - 52.8)/(150.7 - 52.8) = 27.4/97.9 \approx 0.28$: 28~\% $\alpha$, 72~\% $\beta$.
\end{solution}

\begin{solution}{exo:b2:biomolecules:9}
Maltosa mempertahankan satu hemiasetal bebas pada glukosa keduanya, yang
terbuka menjadi aldehida: pereduksi. Dalam sukrosa, kedua \omterm{def:b2:biomolecules:anomer}{karbon anomerik}
terlibat dalam \omterm{def:b2:biomolecules:glycoside}{ikatan glikosidik}: tidak ada hemiasetal, bukan pereduksi.
Asam encer menghidrolisis \omterm{def:b2:biomolecules:glycoside}{ikatan glikosidik}, sebuah asetal: glukosa dan
fruktosa dilepaskan, dan keduanya mereduksi.
\end{solution}

\begin{solution}{exo:b2:biomolecules:10}
Triolein, \qty{885.45}{g/mol}, memerlukan tiga \ce{KOH} (\qty{56.105}{g/mol}): $3 \times 56.105/885.45 =
0.190$~g per g, yaitu bilangan penyabunan 190. Rantai yang lebih pendek
berarti massa molar yang lebih kecil untuk tiga gugus ester yang sama: lebih
banyak mol lemak per gram, lebih banyak KOH.
% ledger: aw:C, aw:H, aw:O, aw:K
\end{solution}

\begin{solution}{exo:b2:biomolecules:11}
Lindungi amina glisina sebagai karbamat \textit{tert}-butoksikarbonil, dan
asam alanina sebagai metil ester. Lakukan kopling dengan karbodiimida.
Lepaskan karbamat dengan asam dan ester dengan basa encer (kondisi
ortogonal, sehingga masing-masing ujung dapat dibebaskan sendiri):
Gly--Ala.
\end{solution}

\begin{solution}{exo:b2:biomolecules:12}
5$'$-ACGCAT-3$'$ (dibaca antiparalel: A berpasangan dengan T, G dengan C).
Pasangan: A--T, T--A, G--C, C--G, G--C, T--A:
$3 \times 2 + 3 \times 3 = 15$ ikatan hidrogen.
\end{solution}

\begin{solution}{pb:b2:biomolecules:1}
\textbf{1.} Asam L-aspartat, L-fenilalanina, dan metanol.
\textbf{2.} Pada keduanya, \ce{COOH} di atas, \omterm{def:b2:biomolecules:amino-acid}{rantai samping} di bawah (\ce{CH2COOH}, \ce{CH2C6H5}), \ce{NH2}
di kiri.
\textbf{3.} (\textit{S}) untuk keduanya.
\textbf{4.} Dua pusat stereogenik: empat stereoisomer.
\textbf{5.} \ce{H2N-CH(CH2COOH)-CO-NH-CH(CH2C6H5)-COOCH3}: \omterm{def:b2:biomolecules:peptide}{ikatan peptida}
menyambungkan kedua residu; esternya adalah \ce{COOCH3} pada ujung-C.
\textbf{6.} Reseptor rasa bersifat kiral: reseptor itu mengikat satu
stereoisomer dan tidak mengikat bayangan cerminnya atau
diastereoisomer-diastereoisomernya, seperti disebutkan pada pembuka bab.
\textbf{7.} \ce{NH3+} ujung-N dan \ce{COOH} \omterm{def:b2:biomolecules:amino-acid}{rantai samping} asam aspartat;
asam ujung-C berupa metil ester, tanpa hidrogen yang asam.
\textbf{8.} Pada pH 2: \ce{NH3+} ($+1$), \ce{COOH} \omterm{def:b2:biomolecules:amino-acid}{rantai samping} (0): $+1$.
\textbf{9.} Pada pH 7: $+1 - 1 = 0$.
\textbf{10.} Pada pH 12: $0 - 1 = -1$.
\textbf{11.} Di antara kedua $\mathrm pK_a$ gugus yang berubah: $(3.9 + 9.9)/2 \approx 6.9$.
\textbf{12.} Dalam dipeptida, karbon amina membawa amida alih-alih
karboksilat: gugus-gugus dan muatan-muatan di sekitarnya berubah, demikian
pula $\mathrm pK_a$-nya (amonium ujung-N sebuah \omterm{def:b2:biomolecules:peptide}{peptida} jelas lebih asam
daripada amonium asam amino bebas). Model itu memberikan bentuk, bukan nilai
yang eksak.
\textbf{13.} Asam dan amina mula-mula membentuk garam; pemanasan akan
menghasilkan campuran (Asp--Asp, rantai yang lebih panjang, gugus
karboksilat yang salah) dan merasemisasi pusat-pusat stereogenik.
\textbf{14.} Gugus aminonya dan gugus karboksilat \omterm{def:b2:biomolecules:amino-acid}{rantai sampingnya}.
\textbf{15.} Pereaksi itu mengubah \ce{OH} asam menjadi gugus pergi yang
baik, sebuah turunan teraktivasi yang diserang amina pada suhu kamar
(substitusi asil).
\textbf{16.} $\beta$-dipeptida, yang tersambung melalui \omterm{def:b2:biomolecules:amino-acid}{rantai samping}:
sebuah isomer aspartam, bukan aspartam.
\textbf{17.} Amina sebagai karbamat benziloksikarbonil dan \omterm{def:b2:biomolecules:amino-acid}{rantai samping}
sebagai benzil ester: keduanya dilepaskan bersama-sama melalui
hidrogenolisis pada paladium, yang membiarkan metil ester tidak tersentuh
(ortogonal terhadapnya).
\textbf{18.} Asam yang teraktivasi mempunyai hidrogen $\alpha$ yang asam;
basa melepaskannya dan pusat stereogeniknya terasemisasi.
\textbf{19.} Lindungi asam aspartat (karbamat pada N, benzil ester pada
\omterm{def:b2:biomolecules:amino-acid}{rantai samping}); aktifkan
$\alpha$-\ce{COOH} yang bebas dengan \omterm{def:b2:biomolecules:peptide}{pereaksi kopling}; tambahkan metil ester
L-fenilalanina; lakukan hidrogenolisis.
\textbf{20.} Metil ester dan \omterm{def:b2:biomolecules:peptide}{ikatan peptida}; ester lebih dahulu, karena amida
jauh kurang reaktif.
\textbf{21.} Dipeptida Asp--Phe (asam bebas) dan metanol.
\textbf{22.} Molekul yang manis itu terpakai: produk-produk hidrolisisnya
tidak menggantikannya.
\textbf{23.} $14 \times 12.011 + 18 \times 1.008 + 2 \times 14.007 + 5 \times 15.999 =
\qty{294.307}{g/mol}$.
\textbf{24.} $0.180/294.307 = \qty{6.116e-4}{mol}$, yaitu \qty{0.612}{mmol}.
\textbf{25.} Satu metanol per aspartam: $0.6116 \times 32.042 = \boldsymbol{\approx \qty{19.6}{mg}}$
metanol.
\end{solution}
